% APÊNDICES: materiais elaborados por você, como questionários e roteiros.
% Para retirar todos os apêndices, comente sua inclusão em estrutura/pos-textual.tex.
\begin{apendicesenv}
\renewcommand{\appendixtocname}{\texorpdfstring{\MakeUppercase{\apendicesname}}{\apendicesname}}

% Comente a linha seguinte para retirar a página de abertura dos apêndices.
\partapendices

% Cada \chapter cria um apêndice, com letra atribuída automaticamente.
% EXEMPLO -- apague este apêndice inteiro ou substitua seu conteúdo.
\chapter{Exemplo de apêndice}
\label{apendice:trabalhos}

Esta página é uma demonstração: o apêndice vem ativado para você ver como
ele aparece. Substitua o conteúdo por um material complementar produzido
na sua pesquisa (questionário, roteiro de entrevista, código-fonte) ou
desative o apêndice comentando a linha correspondente em
\texttt{estrutura/pos-textual.tex}.

\begin{table}[htb]
  \centering
  \ABNTEXfontereduzida
  \caption{Visão geral dos trabalhos selecionados: exemplo fictício}
  \label{tab:appendix-trabalhos-aprovados}
  \begin{tabularx}{\linewidth}{llX}
    \toprule
    \textbf{Identificador} & \textbf{Tipo} & \textbf{Descrição} \\
    \midrule
    S01 & Artigo & Síntese do primeiro trabalho selecionado. \\
    S02 & Dissertação & Síntese do segundo trabalho selecionado. \\
    \bottomrule
  \end{tabularx}
  \fonte{Elaboração própria.}
\end{table}
\end{apendicesenv}
